\item 
\textbf{HydroSense Tech }    \hfill     
\hfill  \textit{Sep 2017 -- Present (Part Time, remote); 
\item  Co-founder / Machine Learning Engineer\hfill May 2025 -- Sep 2025 (Full Time, Singapore)}
\vspace*{-3pt}
\begin{itemize}[leftmargin=0.2in, topsep=1pt]
  \item Co-founded an IoT startup and led ML development, building Python services for large-scale GPU training operations running reliably and efficiently on dedicated GPU clusters with multi-node distributed training using FSDP and DeepSpeed.
  \item Owned training pipeline for large-scale LLM fine-tuning and post-training workflows, configuring and launching multi-node distributed training jobs with FSDP, DeepSpeed, and custom wrappers for optimal throughput.
  \item Contributed to training frameworks including TorchTune, TRL, and Hugging Face Transformers, tuning training parameters, memory footprints, and sharding strategies for optimal performance with PyTorch distributed training.
  \item Worked closely with infrastructure teams to maintain GPU cluster health and utilization using Infiniband, NCCL, Slurm, and Kubernetes, implementing features and fixes for novel use cases in LLM training stack.
  \item \textbf{Patents:} Co-inventor on 3 granted patents covering rainfall sensing hardware, ML-based drift compensation, and intelligent embedded calibration algorithms for industrial monitoring applications.
\end{itemize}

\vspace*{-1pt}
\item \textbf{Seno Medical}\\
\textit{Machine Learning Engineer} \hfill Buffalo, NY \quad \textit{May 2024 -- Aug 2024}
\vspace*{-3pt}
\begin{itemize}[leftmargin=0.2in, topsep=1pt]
  \item Collaborated with domain experts to build a \textbf{2k+ QA} dataset from technical documents; fine-tuned a \textbf{LLaMA-based model (LoRA)} and implemented retrieval-augmented generation (RAG) system with hybrid search for document processing.
  \item Designed an end-to-end \textbf{evaluation harness} with label schema (factuality, coverage, style), automatic metrics using RAGAS and TruLens, and Python error-analysis scripts; reduced hallucination rate from \textbf{30\% to 15\%} and increased \textbf{Recall@5 to 90\%}.
  \item Built and maintained production services using PyTorch with deep familiarity with distributed training via FSDP, DeepSpeed, and DDP, working with large-scale ML systems and training pipelines.
  \item Set up Git-based workflows and production-grade \textbf{observability} using Langfuse and OpenTelemetry, shipping reproducible Docker images with monitoring frameworks for production ML systems with profiling and debugging capabilities.
\end{itemize}

\vspace*{-1pt}
\item \textbf{Roswell Park Comprehensive Cancer Center}\\
\textit{Research Scientist} \hfill Buffalo, NY \quad \textit{May 2023 -- Aug 2023}
\vspace*{-3pt}
\begin{itemize}[leftmargin=0.2in, topsep=1pt]
  \item Designed and implemented a \textbf{multi-modal data pipeline} aligning imaging with structured text reports, including de-identification, schema normalization, and automated QC, to produce clean image–text pairs for generative modeling.
  \item Built \textbf{dataset generation jobs} with chunking strategies, filtering, sampling, and augmentation for embeddings indexing, creating balanced image–report datasets with vector database integration for conditional generation.
  \item Implemented a \textbf{domain-specific latent diffusion model} (Stable Diffusion–style, DDIM/DPMS sampling) trained on \textbf{AWS}, integrating LLM orchestration for conditional image–to–text report generation and synthetic data augmentation.
  \item Automated \textbf{GPU training and experiment tracking} with standardized configs, observability pipelines, and comparison scripts, enabling fast iteration over model variants and dataset settings on distributed AWS infrastructure.
\end{itemize}
